\chapter{Jeux à Somme Nulle}

\section{Préliminaires Topologiques}

\begin{definition}[Semi-continuité]
Une fonction $f:X \to \mathbb{R}$ est \textbf{semi-continue supérieurement (s.c.s.)} si pour tout $\alpha \in \mathbb{R}$, l'ensemble $\{x \in X : f(x) \ge \alpha\}$ est fermé.

Une fonction $f:X \to \mathbb{R}$ est \textbf{semi-continue inférieurement (s.c.i.)} si pour tout $\alpha \in \mathbb{R}$, l'ensemble $\{x \in X : f(x) \le \alpha\}$ est fermé.
\end{definition}

\begin{exercice}[Propriétés à démontrer]
\begin{itemize}
    \item Montrer que l'infimum d'une famille de fonctions s.c.s. est s.c.s.
    \item Montrer que le supremum d'une famille de fonctions s.c.i. est s.c.i.
\end{itemize}
\paragraph{Correction :}
\begin{enumerate}
    \item Soit $f = \inf_{i \in I} f_i$ où chaque $f_i$ est s.c.s.
    $ \{ x : f(x) \ge \alpha \} = \{ x : \forall i, f_i(x) \ge \alpha \} = \bigcap_{i \in I} \{ x : f_i(x) \ge \alpha \} $.
    Chaque ensemble $\{ f_i \ge \alpha \}$ est fermé (déf s.c.s.). Une intersection quelconque de fermés est fermée. Donc $f$ est s.c.s.
    
    \item Soit $g = \sup f_i$ avec $f_i$ s.c.i.
    $ \{ x : g(x) \le \alpha \} = \{ x : \sup f_i(x) \le \alpha \} = \{ x : \forall i, f_i(x) \le \alpha \} = \bigcap \{ f_i \le \alpha \}$.
    Intersection de fermés $\implies$ fermé. Donc $g$ est s.c.i.
\end{enumerate}
\end{exercice}


\begin{exercice}[Caractérisation de la quasi-concavité]
Prouver qu'une fonction $f:X \to \mathbb{R}$ est quasi-concave si et seulement si, pour tout $\alpha \in \mathbb{R}$, l'ensemble de niveau supérieur $L_{\alpha} = \{x \in X : f(x) \ge \alpha\}$ est un ensemble convexe.
\end{exercice}

\paragraph{Correction :}
\begin{itemize}
    \item $\Rightarrow$ : Supposons $f$ quasi-concave. Par définition $f(\lambda x + (1-\lambda)y) \ge \min(f(x), f(y))$.
    Soient $x, y \in L_\alpha$. Alors $f(x) \ge \alpha$ et $f(y) \ge \alpha$. Donc $\min \ge \alpha$.
    D'où $f(z) \ge \alpha$ pour tout $z$ segment $[x,y]$. Donc $z \in L_\alpha$. $L_\alpha$ est convexe.
    \item $\Leftarrow$ : Supposons les $L_\alpha$ convexes.
    Soit $\alpha = \min(f(x), f(y))$. Alors $x \in L_\alpha$ et $y \in L_\alpha$.
    Par convexité, le segment joint est dans $L_\alpha$, donc $f(\lambda x + \dots) \ge \alpha$, c.q.f.d.
\end{itemize}

\section{Définitions et Analogies}

\begin{definition}[Jeu à somme nulle à 2 joueurs]
Un jeu à somme nulle est un triplet $(X, Y, u)$ où :
\begin{itemize}
    \item $X$ est l'ensemble des stratégies du Joueur 1 (maximizant).
    \item $Y$ est l'ensemble des stratégies du Joueur 2 (minimizant).
    \item $u: X \times Y \rightarrow \mathbb{R}$ représente le gain du Joueur 1.
    \item Le gain du Joueur 2 est \textbf{identiquement opposé} : $-u(x, y)$.
\end{itemize}
Ce modèle capture l'essence du \textbf{conflit pur} : ce que l'un gagne, l'autre le perd nécessairement.
\end{definition}

\paragraph{Analogies Sportives et Classiques :}
\begin{itemize}
    \item \textbf{Penalty au Football :} Si le tireur marque (+1), le gardien encaisse (-1). C'est l'archétype du jeu à somme nulle.
    \item \textbf{Matching Pennies (Pile ou Face) :} Chaque joueur choisit H ou T. Si les choix correspondent, J1 gagne 1 (J2 perd 1). Sinon, J1 perd 1 (J2 gagne 1).
    $$M = \begin{pmatrix} 1 & -1 \\ -1 & 1 \end{pmatrix}$$
\end{itemize}

\section{Valeurs Maximin et Minimax (Intuition)}

Soit $v_1(x) = \inf_{y \in Y} u(x, y)$ le pire résultat possible pour J1 s'il joue $x$.
Le joueur 1, prudent, cherchera à maximiser cette sécurité :
$$\underline{v} = \sup_{x \in X} \inf_{y \in Y} u(x, y) \quad \text{(Valeur Maximin ou Niveau de Sécurité)}$$

De même, le joueur 2 veut minimiser le gain de J1 (donc maximiser le sien). Il anticipe le meilleur coup de J1 :
$$\overline{v} = \inf_{y \in Y} \sup_{x \in X} u(x, y) \quad \text{(Valeur Minimax ou Plafond de Paiement)}$$

\begin{definition}[Valeur du Jeu]
Un jeu à somme nulle $G=(X, Y, u)$ possède une \textbf{valeur} $V \in \mathbb{R}$ si :
$$ V = \underline{v} = \overline{v} $$
C'est-à-dire si le maximin est égal au minimax.
\end{definition}

\begin{remarque}[Valeur de Sécurité]
La valeur minimax $\overline{v}$ peut être vue comme une sorte de "valeur de sécurité" pour le joueur 2 (dans le sens où le joueur 2 est sûr de ne pas avoir à payer plus qu'une valeur très proche de $\overline{v}$).
\end{remarque}

\begin{proposition}[Inégalité Fondamentale]
Pour toute fonction $u$, on a toujours :
$$\underline{v} \le \overline{v}$$
L'inverse est faux en général (sauf s'il y a un point-selle).
\end{proposition}

\section{Outils Topologiques (Préalables)}
Le théorème de Sion requiert des notions précises de continuité et de convexité.

\begin{definition}[Semi-continuité]
Soit $f: X \to \mathbb{R}$.
\begin{itemize}
    \item $f$ est \textbf{semi-continue supérieurement (s.c.s.)} si pour tout $\lambda$, l'ensemble de niveau supérieur $\{x : f(x) \ge \lambda\}$ est \textbf{fermé}.
    \item $f$ est \textbf{semi-continue inférieurement (s.c.i.)} si pour tout $\lambda$, l'ensemble de niveau inférieur $\{x : f(x) \le \lambda\}$ est \textbf{fermé}.
\end{itemize}
\end{definition}

\begin{definition}[Quasi-convexité]
$f$ est \textbf{quasi-convexe} si ses ensembles de niveau inférieur $\{x : f(x) \le \lambda\}$ sont \textbf{convexes}. (L'opposé d'une fonction quasi-concave).
\end{definition}

\subsection{Le Lemme d'Intersection}
Ce lemme est crucial pour la preuve de Sion :
\begin{lemme}[de Séparation (Préliminaire)]
Soient $A$ et $B$ deux ensembles convexes compacts non vides et disjoints. Soit $H$ l'hyperplan qui les sépare strictement (Théorème de Hahn-Banach). Si un ensemble convexe $C$ a une intersection non vide avec $A$ et avec $B$, alors $C \cap H \neq \emptyset$.
\end{lemme}

\begin{lemme}[d'Intersection]
Soient $C_1, \dots, C_n$ des ensembles convexes compacts non vides. Si $\cup_{i=1}^n C_i$ est convexe et si pour tout $j \in \{1,\dots,n\}$, $\cap_{i \neq j} C_i \neq \emptyset$, alors $\cap_{i=1}^n C_i \neq \emptyset$.
\end{lemme}

\begin{proof}[Démonstration (par récurrence sur n)]
\begin{itemize}
    \item \textbf{Initialisation ($n=2$) :} Si $C_1 \cup C_2$ est convexe avec $C_1 \cap C_2 = \emptyset$, alors par Hahn-Banach il existe un hyperplan $H$ séparant strictement $C_1$ et $C_2$. Comme $C_1 \cup C_2$ est convexe et intersecte les deux côtés de $H$, il doit intersecter $H$. Or $H$ ne rencontre ni $C_1$ ni $C_2$, contradiction. Donc $C_1 \cap C_2 \neq \emptyset$.
    \item \textbf{Hérédité :} On suppose le résultat vrai pour $n-1$. (La suite de la preuve utilise la construction d'intersections de type $C_i \cap C_n$).
\end{itemize}
\end{proof}

\subsection{Notation Matricielle de Von Neumann}
Pour un jeu fini défini par une matrice $A$, on note le gain :
\[ val(A) = \max_{x \in \Delta(I)} \min_{y \in \Delta(J)} xAy = \min_{y \in \Delta(J)} \max_{x \in \Delta(I)} xAy \]
où $x$ est un vecteur ligne et $y$ un vecteur colonne.
\section{Théorème du Minimax de von Neumann}

Ce théorème établit une connexion profonde entre l'équilibre de Nash et les valeurs maximin/minimax dans les jeux à somme nulle.

\begin{theoreme}[Minimax de von Neumann]
Pour un jeu à deux joueurs à somme nulle, les deux affirmations suivantes sont équivalentes :
\begin{enumerate}
    \item Le jeu possède un \textbf{équilibre de Nash} $(x^*, y^*)$.
    \item Le jeu a une \textbf{valeur}, c'est-à-dire $\underline{v} = \overline{v} = V$, et chaque joueur dispose d'une stratégie optimale.
\end{enumerate}
\textbf{Conséquence :} Si ces conditions sont remplies, le gain à l'équilibre est $(V, -V)$.
\end{theoreme}
\section{Équilibre $\epsilon$-Nash et Valeur}

\begin{definition}[$\epsilon$-équilibre]
Un profil $(x, y)$ est un $\epsilon$-équilibre si personne ne peut gagner plus de $\epsilon$ en déviant :
$$u(d_x, y) - \epsilon \le u(x, y) \quad \text{et} \quad -u(x, d_y) - \epsilon \le -u(x, y)$$
\end{definition}

\begin{proposition}
Un jeu possède une valeur si et seulement si pour tout $\epsilon > 0$, il existe un $\epsilon$-équilibre.
\end{proposition}

\section{Étude Détaillée : Théorème de Minimax de Sion}

\subsection{Hypothèses et Définitions}
Soit un jeu $(X, Y, u)$ où:
\begin{itemize}
    \item $X$ est un ensemble compact et convexe.
    \item $Y$ est un ensemble convexe.
    \item $u$ est quasi-concave et semi-continue supérieurement (s.c.s.) en $x$.
    \item $u$ est quasi-convexe et semi-continue inférieurement (s.c.i.) en $y$.
\end{itemize}

\subsection{Partie 1 : Construction de la Contradiction}

\paragraph{Question 1} Expliquer pourquoi il existe un réel $v$ tel que :
$$ \sup_{x \in X} \inf_{y \in Y} u(x,y) < v < \inf_{y \in Y} \sup_{x \in X} u(x,y) $$
\textbf{Correction :} On suppose par l'absurde que le jeu n'a pas de valeur, soit $\underline{v} < \overline{v}$. Par la propriété de densité des réels, on peut choisir un $v$ entre ces deux bornes. Si les valeurs sont finies, on peut prendre $v = \frac{\underline{v} + \overline{v}}{2}$. Si $\overline{v} = +\infty$, on peut prendre $v = \underline{v} + 1$.

\paragraph{Question 2} Expliquer pourquoi (1) $\forall x \in X, \exists y \in Y : u(x,y) < v$ et (2) $\forall y \in Y, \exists x \in X : u(x,y) > v$.
\textbf{Correction :} 
(1) Puisque $\sup_x (\inf_y u(x,y)) < v$, alors pour n'importe quel $x$, la valeur minimale $\inf_y u(x,y)$ est strictement inférieure à $v$. Il existe donc forcément un $y$ (qui dépend de $x$) tel que $u(x,y) < v$. 
(2) De même, $\inf_y (\sup_x u(x,y)) > v$ implique que pour tout $y$, le meilleur gain possible pour le joueur 1 est au-dessus de $v$, donc il existe un $x$ tel que $u(x,y) > v$.

\paragraph{Question 3} Pour chaque $y \in Y$ et $\epsilon > 0$, soit $X_y^\epsilon = \{x \in X : u(x,y) < v - \epsilon\}$. Prouver que c'est un ouvert.
\textbf{Correction :} Le complémentaire de cet ensemble dans $X$ est $\{x \in X : u(x,y) \geq v - \epsilon\}$. Comme l'application $x \mapsto u(x,y)$ est semi-continue supérieurement (s.c.s.), ses ensembles de niveau supérieur sont fermés par définition. Le complémentaire $X_y^\epsilon$ est donc un ouvert.

\paragraph{Question 4} Prouver que pour tout $x \in X$, il existe $y \in Y$ et $\epsilon > 0$ tels que $u(x,y) < v - \epsilon$.
\textbf{Correction :} D'après la question 2, pour tout $x$, il existe $y$ tel que $u(x,y) < v$. En posant $\epsilon = \frac{v - u(x,y)}{2} > 0$, on obtient bien $u(x,y) = v - 2\epsilon < v - \epsilon$.

\paragraph{Question 5} Prouver que $X \subset \bigcup_{y \in Y, \epsilon > 0} X_y^\epsilon$.
\textbf{Correction :} Puisque chaque point $x$ de $X$ appartient à au moins un ensemble $X_y^\epsilon$ (d'après la Question 4), l'union de tous ces ouverts recouvre nécessairement $X$.

\paragraph{Question 6} Prouver qu'il existe un ensemble fini $Y_0 \subset Y$ et $\epsilon > 0$ tels que $X \subset \bigcup_{y \in Y_0} X_y^\epsilon$.
\textbf{Correction :} L'ensemble $X$ est compact. Par la propriété de Borel-Lebesgue, de tout recouvrement ouvert, on peut extraire un sous-recouvrement fini. Il existe donc $y_1, \dots, y_n$ et $\epsilon_i$ associés. En prenant $\epsilon = \min(\epsilon_1, \dots, \epsilon_n)$, on conserve l'inclusion car $X_{y_i}^{\epsilon_i} \subset X_{y_i}^{\epsilon}$ pour $\epsilon \leq \epsilon_i$.

\paragraph{Question 7} Prouver que $Y' := \text{conv}(Y_0)$ est compact et convexe.
\textbf{Correction :} $Y'$ est l'enveloppe convexe d'un nombre fini de points. Il est convexe par définition. Comme on travaille dans un espace de dimension finie (le span de $Y_0$), il est borné. Pour la compacité, toute suite dans $Y'$ a ses coefficients de mélange dans le simplexe (qui est compact), donc par Bolzano-Weierstrass, on peut extraire une sous-suite convergente dont la limite est encore dans $Y'$.

\paragraph{Question 8} Prouver que : $\sup_{x \in X} \inf_{y \in Y_0} u(x,y) < v < \inf_{y \in Y'} \sup_{x \in X} u(x,y)$.
\textbf{Correction :}
(Gauche) Puisque $X$ est couvert par les $X_y^\epsilon$ pour $y \in Y_0$, pour tout $x$, il existe un $y \in Y_0$ tel que $u(x,y) < v - \epsilon$. L'infimum sur $Y_0$ est donc toujours inférieur à $v$.
(Droite) Comme $Y' \subset Y$, l'infimum sur $Y'$ est supérieur ou égal à l'infimum sur $Y$, qui est lui-même strictement supérieur à $v$ d'après l'hypothèse de la Question 1.

\subsection{Partie 2 : Symétrisation}

\paragraph{Question 11} Prouver qu'il existe un ensemble fini $X_0 \subset X$ tel que $Y' \subset \bigcup_{x \in X_0} Y_x^\epsilon$.
\textbf{Correction :} On applique le même raisonnement que précédemment : $Y'$ est compact et $u$ est s.c.i. en $y$, donc les ensembles $Y_x^\epsilon = \{y \in Y' : u(x,y) > v + \epsilon\}$ sont ouverts. Par compacité de $Y'$, on extrait un sous-recouvrement fini $X_0$.

\paragraph{Question 13} Prouver que $\sup_{x \in X'} \inf_{y \in Y'} u(x,y) < v < \inf_{y \in Y'} \sup_{x \in X_0} u(x,y)$ (où $X' = \text{conv}(X_0)$).
\textbf{Correction :} Cela découle de la monotonie du sup et de l'inf sur les sous-ensembles compacts $X'$ et $Y'$.

\paragraph{Question 16} Justifier l'existence de $X_0$ et $Y_0$ minimaux satisfaisant ces inégalités.
\textbf{Correction :} Si un élément de $X_0$ ou $Y_0$ n'est pas nécessaire pour maintenir l'inégalité, on peut simplement le retirer de l'ensemble fini jusqu'à ce que l'ensemble soit minimal pour l'inclusion.

\subsection{Partie 3 : Lemme d'Intersection et Contradiction}

\paragraph{Question 17} Pour $x \in X_0$, soit $A_x = \{y \in Y' : u(x,y) \leq v\}$. Prouver que $A_x$ est compact et convexe.
\textbf{Correction :} 
\begin{itemize}
    \item \textbf{Compact :} $u$ est s.c.i. en $y$, donc $A_x$ est fermé. Étant contenu dans le compact $Y'$, il est compact.
    \item \textbf{Convexe :} $u$ est quasi-convexe en $y$, donc ses ensembles de niveau inférieur sont convexes par définition.
\end{itemize}

\paragraph{Question 18} Prouver que $\bigcap_{x \in X_0} A_x = \emptyset$.
\textbf{Correction :} Si l'intersection contenait un point $y^*$, on aurait $u(x, y^*) \leq v$ pour tout $x \in X_0$. Cela impliquerait $\sup_{x \in X_0} u(x, y^*) \leq v$, ce qui contredit l'inégalité de droite de la Question 13 : $\inf_{y \in Y'} \sup_{x \in X_0} u(x, y) > v$.

\paragraph{Question 20} Prouver que $\bigcup_{x \in X_0} A_x$ n'est pas convexe et n'est pas égal à $Y'$.
\textbf{Correction :} Par le Lemme d'Intersection, si une famille d'ensembles compacts convexes $A_x$ a une intersection totale vide mais que toutes ses sous-familles ont une intersection non vide (par minimalité de $X_0$), alors leur union n'est pas convexe. En particulier, elle ne peut pas être égale à l'ensemble convexe $Y'$.

\paragraph{Question 21 \& 22} Prouver l'existence de $y_0 \in Y'$ tel que $u(x, y_0) > v$ pour tout $x \in X'$.
\textbf{Correction :} Comme l'union des $A_x$ ne couvre pas $Y'$, il existe $y_0 \in Y' \setminus \bigcup A_x$. Donc $y_0$ n'appartient à aucun $A_x$, d'où $\forall x \in X_0, u(x, y_0) > v$. Par quasi-concavité de $u$ en $x$, cette propriété reste vraie pour toute combinaison convexe $x \in X' = \text{conv}(X_0)$.

\paragraph{Question 23} Finir la preuve.
\textbf{Correction :} Par un argument symétrique, il existe $x_0 \in X'$ tel que pour tout $y \in Y'$, $u(x_0, y) < v$. En évaluant au point $(x_0, y_0)$, on obtient $u(x_0, y_0) > v$ et $u(x_0, y_0) < v$. C'est une contradiction. Donc $\underline{v} = \overline{v}$.

\paragraph{Question 25} Prouver l'existence d'une stratégie optimale.
\textbf{Correction :} L'application $f(x) = \inf_{y \in Y} u(x,y)$ est l'infimum d'une famille de fonctions s.c.s., elle est donc s.c.s.. D'après le théorème de Weierstrass, une fonction s.c.s. sur un compact ($X$) atteint son maximum. Il existe donc un $\overline{x}$ tel que $f(\overline{x}) = \sup_x f(x)$, ce qui est la stratégie optimale.

\begin{exercice}[Contre-exemple de Sion-Wolfe]
Soit le jeu sur $S_1 = S_2 = [0,1]$ avec :
$$u_1(x,y) = \begin{cases} -1 & \text{si } x < y < x + 1/2 \\ 0 & \text{si } x = y \text{ ou } y = x + 1/2 \\ 1 & \text{sinon} \end{cases}$$
Montrer que ce jeu ne possède pas de valeur (le Min-Max diffère du Max-Min) car les conditions de régularité (semi-continuité) de Sion ne sont pas vérifiées.
\end{exercice}

\paragraph{Correction :}
\begin{itemize}
    \item \textbf{Calcul du Maximin ($\underline{v}$)} :
    Pour $x$ fixé, le joueur 2 peut choisir $y=x + \epsilon$ (avec $\epsilon$ petit). Alors $x < y < x+1/2$, donc $u=-1$.
    Ainsi $\inf_y u(x,y) = -1$.
    Et $\sup_x (-1) = -1$. Donc $\underline{v} = -1$.
    
    \item \textbf{Calcul du Minimax ($\overline{v}$)} :
    Pour $y$ fixé, le joueur 1 peut choisir $x = y$ ou très proche.
    Si J1 choisit $x=y$, $u=0$.
    Mais J1 peut choisir $x$ tel que $y$ ne soit pas dans $]x, x+1/2[$. Par exemple $x$ tel que $x > y$. Alors $u=1$.
    Ainsi $\sup_x u(x,y) = 1$.
    Et $\inf_y (1) = 1$. Donc $\overline{v} = 1$.
    
    \item \textbf{Conclusion :} $-1 \ne 1$. Le jeu n'a pas de valeur. Cela vient de la discontinuité grossière des gains.
\end{itemize}

\section{Exercices Complémentaires}

\subsection{Exercice Matriciel 4x4}
Calculer $\underline{v}$ et $\overline{v}$ pour :
$$M = \begin{pmatrix} -1 & 2 & -1 & 0 \\ 1 & 3 & 3 & 5 \\ -9 & -2 & 8 & 2 \\ -1 & 4 & -4 & 5 \end{pmatrix}$$
\textit{Méthode :}
\begin{itemize}
    \item Maximin (Ligne) : Min de chaque ligne $\to (-1, 1, -9, -4)$. Le max est $1$ (Ligne 2). Donc $\underline{v} = 1$.
    \item Minimax (Colonne) : Max de chaque colonne $\to (1, 4, 8, 5)$. Le min est $1$ (Colonne 1). Donc $\overline{v} = 1$.
    \item Conclusion : Point selle $(L2, C1)$, valeur $v=1$.
\end{itemize}

\subsection{Exercice "No Value" (Cours)}
Considérons le jeu à somme nulle donné par la matrice suivante (actions J1: $i$, actions J2: $j$) :
$$ A = \begin{pmatrix} 1 & 2 & -1 & 0 \\ 0 & 1 & 3 & 5 \\ 9 & -2 & 8 & 2 \\ 1 & 4 & -4 & 5 \end{pmatrix} $$

\textbf{Analyse des gains :}
Par exemple, si $i=1$ et $j=2$ (J1 joue 1, J2 joue 2), alors J1 gagne 2 et J2 gagne -2.

\textbf{Calcul des valeurs pures :}
\begin{itemize}
    \item \textbf{Maximin ($\underline{V}$)} :
    \begin{itemize}
        \item Ligne 1 : min = -1
        \item Ligne 2 : min = 0
        \item Ligne 3 : min = -2
        \item Ligne 4 : min = -4
    \end{itemize}
    Le maximum de ces minimums est $\underline{V} = 0$.

    \item \textbf{Minimax ($\overline{V}$)} :
    \begin{itemize}
        \item Col 1 : max = 9
        \item Col 2 : max = 4
        \item Col 3 : max = 8
        \item Col 4 : max = 5
    \end{itemize}
    Le minimum de ces maximums est $\overline{V} = 4$.
\end{itemize}

\textbf{Conclusion :}
$$ \underline{V} = 0 < 4 = \overline{V} $$
Le jeu n'a \textbf{pas de valeur} (en stratégies pures). Il n'y a donc pas d'équilibre de Nash ni de consensus possible.

\subsection{Jeu 2x2 sans valeur pure (Exemple du cours)}
$$M = \begin{pmatrix} 2 & 1 \\ 0 & 2 \end{pmatrix}$$
\begin{enumerate}
    \item \textbf{Valeurs Pures :}
    \begin{itemize}
        \item Maximin (J1) : min ligne 1 = 1, min ligne 2 = 0. Max = 1. Donc $\underline{v} = 1$.
        \item Minimax (J2) : max col 1 = 2, max col 2 = 2. Min = 2. Donc $\overline{v} = 2$.
    \end{itemize}
    Comme $\underline{v} < \overline{v}$, il n'y a pas de point selle (pas de consensus).

    \item \textbf{Stratégies Mixtes (Méthode de l'indifférence) :}
    Soient $p$ la probabilité que J1 joue Haut, et $q$ la probabilité que J2 joue Gauche.
    \begin{itemize}
        \item \textbf{Indifférence de J1 :} J1 doit gagner la même chose si J2 joue G ou D (pour que J2 soit prêt à mixer).
         \textit{Note : L'image des notes inverse parfois la logique, mais théoriquement :}
         J1 rend J2 indifférent :
         $$2p + 0(1-p) = 1p + 2(1-p)$$
         $$2p = p + 2 - 2p \implies 3p = 2 \implies p = 2/3$$
        
        \item \textbf{Indifférence de J2 :} J2 rend J1 indifférent :
        $$2q + 1(1-q) = 0q + 2(1-q)$$
        $$q + 1 = 2 - 2q \implies 3q = 1 \implies q = 1/3$$
    \end{itemize}
    \textbf{Équilibre de Nash :} $((2/3, 1/3), (1/3, 2/3))$.
    \textbf{Valeur du jeu :} Calculons le gain de J1 :
    $u_1 = 2(1/3) + 1(2/3) = 4/3$ (si J1 joue H contre q).
    Vérif: $0(1/3) + 2(2/3) = 4/3$. OK.
\end{enumerate}

\begin{exercice}[Lien entre Valeur et $\epsilon$-Nash]
Prouver qu'un jeu à somme nulle possède une valeur $V$ si et seulement si, pour tout $\epsilon > 0$, il existe un profil de stratégies $(x^*_\epsilon, y^*_\epsilon)$ tel que :
$$ \underline{v}(x^*_\epsilon) \ge V - \epsilon \quad \text{et} \quad \overline{v}(y^*_\epsilon) \le V + \epsilon $$
\end{exercice}

\paragraph{Correction :}
\begin{itemize}
    \item $\Rightarrow$ Si $V$ existe, $\underline{v} = \overline{v} = V$. Il suffit de prendre $x^*_\epsilon$ une stratégie $\epsilon$-optimale pour le maximin, et $y^*_\epsilon$ pour le minimax.
    \item $\Leftarrow$ On a toujours $\underline{v} \le \overline{v}$.
    L'hypothèse donne : $V - \epsilon \le \underline{v}(x^*_\epsilon) \le \underline{v} \le \overline{v} \le \overline{v}(y^*_\epsilon) \le V + \epsilon$.
    En faisant tendre $\epsilon \to 0$, on obtient $V \le \underline{v} \le \overline{v} \le V$, donc égalité partout.
\end{itemize}
